% =====================================================================
%  Capitolo 4 — Metodologia sperimentale
% =====================================================================
\chapter{Metodologia sperimentale}
\label{cap:metodologia}

Questo capitolo traduce le domande di ricerca in un disegno sperimentale: ipotesi, dati,
protocolli di split, configurazioni di addestramento e di valutazione, implementazione, budget
computazionale e analisi statistica. Il principio guida è che ogni confronto cambi \emph{un solo}
fattore alla volta e che ogni differenza venga riportata con la sua variabilità.

% ---------------------------------------------------------------------
\section{Ipotesi e disegno sperimentale}
\label{sec:disegno}

Le domande di ricerca della \cref{sec:domande} si traducono in ipotesi verificabili. Nel seguito
$\Delta$ indica la differenza di mIoU di test tra la variante e il suo riferimento, mediata sui
seed. Ogni ipotesi è formulata in modo che anche il suo rifiuto sia un risultato informativo.

\begin{ipotesi}[\RQ{1}]\label{ip:lam}
A parità di encoder e protocollo, le connessioni di salto con attenzione lineare migliorano la
mIoU rispetto alle connessioni identità: $\Delta_{\mathrm{LAM}}>0$.
\end{ipotesi}
\begin{ipotesi}[\RQ{2}]\label{ip:vq}
Un codebook sano (perplexity alta durante tutto l'addestramento) non peggiora la mIoU:
$\Delta_{\mathrm{VQ}}\ge0$. Con le sole connessioni di salto libere la perplexity diminuisce
durante l'addestramento (\cref{prop:aggiramento}).
\end{ipotesi}
\begin{ipotesi}[\RQ{3}]\label{ip:dino}
Il vantaggio relativo del pre-addestramento DINO su ImageNet cresce al diminuire della frazione
di etichette: $\Delta_{\mathrm{DINO}}(1\%)>\Delta_{\mathrm{DINO}}(100\%)$.
\end{ipotesi}
\begin{ipotesi}[\RQ{4}]\label{ip:aug}
La variante \enquote{aerea} delle augmentation produce un encoder migliore della variante
standard, soprattutto con poche etichette.
\end{ipotesi}
\begin{ipotesi}[\RQ{6}]\label{ip:ndsm}
L'nDSM come quarto canale migliora la mIoU, con i guadagni maggiori sulle classi edifici e alberi.
\end{ipotesi}

La \RQ{5} è descrittiva: misura costi e accuratezze senza un'ipotesi direzionale. La
\cref{tab:disegno} collega ogni domanda alla variabile manipolata, al riferimento e alla
campagna di esperimenti che la esegue.

\begin{table}[t]
  \centering
  \small
  \caption[Disegno sperimentale per domanda di ricerca]{Disegno sperimentale: per ogni domanda di
    ricerca, il fattore che varia, il riferimento, la campagna che la esegue e il numero di
    passi (pre-addestramenti, training, benchmark) con tre seed.}
  \label{tab:disegno}
  \begin{tabular}{|c|p{3.5cm}|p{2.7cm}|l|c|}
    \hline
    \textbf{RQ} & \textbf{Fattore} & \textbf{Riferimento} & \textbf{Campagna} & \textbf{Passi} \\
    \hline
    \RQ{1} & connessioni di salto: LAM / identità & \maresunet{} R18 ImageNet & \codice{rq1\_attention} & 6 \\
    \hline
    \RQ{2} & neck: identità / VQ; encoder R18/R34; pre-addestr. ImageNet/DINO & stessa configurazione senza VQ & \codice{thesis\_grid} & 26 \\
    \hline
    \RQ{3} & frazione di etichette $\{1, 10, 25, 100\}\%$ $\times$ pre-addestr. & ImageNet, stessa frazione & \codice{rq3\_low\_label} & 38 \\
    \hline
    \RQ{4} & augmentation DINO: standard / aerea & DINO standard & \codice{rq4\_aerial} & 14 \\
    \hline
    \RQ{5} & encoder R18/R34/R50, Swin-T, ViT & \maresunet{} R18 & \makecell[l]{\codice{rq5\_efficiency}\\\codice{transformers}} & 10 + 10 \\
    \hline
    \RQ{6} & ingresso RGB / RGB + nDSM (ufficiale, calcolato) & RGB & \makecell[l]{\codice{rq6\_ndsm}\\\codice{vaihingen}} & 9 + 9 \\
    \hline
    -- & protocollo del 2022 & tab.~\ref{tab:risultati-2022} & \codice{thesis2022} & 3 \\
    \hline
  \end{tabular}
\end{table}

% ---------------------------------------------------------------------
\section{Dati}
\label{sec:dati}

\subsection{ISPRS Potsdam e Vaihingen}

Il dataset principale è ISPRS Potsdam, il secondo benchmark è ISPRS
Vaihingen~\cite{rottensteiner2012isprs,isprs2dlabeling} (\cref{tab:dataset}). Le classi e la
palette delle maschere sono le stesse per i due dataset (\cref{tab:classi}). Rispetto al codice
del 2022, l'ordine degli indici segue la convenzione ISPRS/MMSegmentation, con il clutter per
ultimo: così il valore ignorato (255) non può coincidere per errore con una classe reale.

\begin{table}[t]
  \centering
  \small
  \caption[Caratteristiche dei dataset ISPRS]{Caratteristiche dei dataset ISPRS usati.}
  \label{tab:dataset}
  \begin{tabular}{|l|p{4.6cm}|p{4.6cm}|}
    \hline
    & \textbf{Potsdam} & \textbf{Vaihingen} \\
    \hline
    Aree & 38 tile $6000\times6000$ & 33 aree, circa $1900\times2500$ \\
    GSD & \SI{5}{\centi\meter} & \SI{9}{\centi\meter} \\
    Immagine & RGB (\codice{2\_Ortho\_RGB}) & IR-R-G in falso colore (\codice{top}) \\
    Altezze & DSM assoluto e nDSM ufficiale & solo DSM assoluto (nDSM calcolato) \\
    Etichette & 6 classi, piene ed \emph{eroded} & 6 classi, piene ed \emph{eroded} \\
    Metrica principale & mIoU con tutte le classi & mIoU senza clutter (quasi assente) \\
    \hline
  \end{tabular}
\end{table}

\begin{table}[t]
  \centering
  \small
  \caption[Classi ISPRS e palette delle maschere]{Classi ISPRS, indici usati in questo lavoro e
    colori delle maschere.}
  \label{tab:classi}
  \begin{tabular}{|c|l|c|c|}
    \hline
    \textbf{Indice} & \textbf{Classe} & \textbf{RGB} & \textbf{Colore} \\
    \hline
    0 & Superfici impermeabili & (255, 255, 255) & \ClasseBoxX{ClsImpervious} \\
    1 & Edifici               & (0, 0, 255)     & \ClasseBoxX{ClsBuilding} \\
    2 & Vegetazione bassa     & (0, 255, 255)   & \ClasseBoxX{ClsLowVeg} \\
    3 & Alberi                & (0, 255, 0)     & \ClasseBoxX{ClsTree} \\
    4 & Auto                  & (255, 255, 0)   & \ClasseBoxX{ClsCar} \\
    5 & Clutter / sfondo      & (255, 0, 0)     & \ClasseBoxX{ClsClutter} \\
    255 & non etichettato (ignorato) & -- & -- \\
    \hline
  \end{tabular}
\end{table}

\subsection{Problemi noti dei file}

L'ispezione dei file distribuiti ha rivelato alcuni difetti che, se ignorati, falserebbero
silenziosamente i risultati (\cref{tab:difetti-file}). La preparazione dei dati li gestisce con
regole esplicite: un ordine di priorità tra le cartelle delle etichette, il rifiuto delle
etichette con più dello 0,5\% di colori fuori palette e un controllo delle dimensioni con
tolleranza di 2 pixel.

\begin{table}[t]
  \centering
  \small
  \caption[Difetti noti dei file ISPRS]{Difetti noti dei file ISPRS e loro gestione.}
  \label{tab:difetti-file}
  \begin{tabular}{|p{4.8cm}|p{2.9cm}|p{4.2cm}|}
    \hline
    \textbf{Difetto} & \textbf{Effetto se ignorato} & \textbf{Gestione} \\
    \hline
    Etichetta di \tile{4\_12} in \codice{5\_Labels\_all} corrotta (nessun colore ISPRS) &
      una tile di validation tutta ignorata & si usa la copia in \codice{5\_Labels\_for\_\allowbreak participants} \\
    \hline
    Etichetta di \tile{6\_7}: lo 0,7\% dei pixel \enquote{auto} ha colore (252, 255, 0) &
      pixel ignorati & stessa priorità tra cartelle \\
    \hline
    DSM e nDSM di \tile{3\_13} larghi 5999 pixel & disallineamento con l'immagine & una colonna replicata dal bordo \\
    \hline
    nDSM \codice{ownapproach}: 11 file su 38 vuoti & ablation impossibile & si usa la variante \codice{lastools} \\
    \hline
  \end{tabular}
\end{table}

% ---------------------------------------------------------------------
\section{Protocolli di split}
\label{sec:split}

Gli split sono sempre \textbf{per tile intere}: le tile vengono assegnate agli split \emph{prima}
del taglio in ritagli, quindi nessun ritaglio può contenere pixel di due split e nessuna porzione
di scena compare sia in training sia in valutazione~\cite{kaufman2012leakage}
(\cref{fig:split}).

\begin{description}[style=nextline, leftmargin=1.2em]
  \item[Potsdam \enquote{standard}.] È lo split di riferimento della
  letteratura~\cite{mmseg2020}: 24 tile di training, che coincidono con quelle distribuite con le
  etichette ai partecipanti del contest, e 14 tile di test. La validation usa due tile di
  training (\tile{4\_12} e \tile{7\_11}) tenute fuori per intero. Su questo protocollo si basano
  tutte le conclusioni.
  \item[Potsdam \enquote{tesi 2022}.] Riproduce il protocollo del 2022: esclude le stesse dieci
  aree, usa 22 tile di training e 6 di validation. Il testo originale non specifica quali sei:
  qui sono fissate a \tile{3\_14}, \tile{4\_11}, \tile{4\_15}, \tile{5\_13}, \tile{6\_15} e
  \tile{7\_11}. Il \enquote{test} coincide con la validation, valutata su tile intere. Serve solo
  al confronto con la \cref{tab:risultati-2022}.
  \item[Vaihingen \enquote{standard}.] 16 aree di training e 17 di test. La validation usa le
  aree 11 e 30. L'area 26 resta in training perché contiene gran parte del poco clutter
  disponibile.
\end{description}

\begin{figure}[t]
  \centering
  \resizebox{\linewidth}{!}{% Mappa delle 38 tile di Potsdam (righe 2-7, colonne 7-15) con l'assegnazione agli split:
% (a) protocollo standard della letteratura, (b) protocollo della tesi del 2022.
\begin{tikzpicture}[x=7.2mm, y=7.2mm]
  \newcommand{\cella}[4]{% riga, colonna, colore di riempimento, colore del testo
    \draw[black!45, line width=0.3pt, fill=#3] (#2-0.5,-#1-0.5) rectangle (#2+0.5,-#1+0.5);
    \node[font=\tiny, text=#4] at (#2,-#1) {#1\_#2};}
  \newcommand{\esclusa}[2]{%
    \draw[black!45, line width=0.3pt, pattern=north east lines, pattern color=black!35]
      (#2-0.5,-#1-0.5) rectangle (#2+0.5,-#1+0.5);
    \node[font=\tiny, text=black!60, fill=white, inner sep=0.5pt] at (#2,-#1) {#1\_#2};}
  \newcommand{\assi}{%
    \foreach \c in {7,...,15} \node[font=\tiny, text=SapGrigio] at (\c,-1.2) {\c};
    \foreach \r in {2,...,7} \node[font=\tiny, text=SapGrigio] at (6.0,-\r) {\r};}

  % ---------------- (a) standard ------------------------------------------------
  \begin{scope}
    \assi
    \foreach \r/\c in {2/10,2/11,2/12,3/10,3/11,3/12,4/10,4/11,5/10,5/11,5/12,
                       6/7,6/8,6/9,6/10,6/11,6/12,7/7,7/8,7/9,7/10,7/12}
      {\cella{\r}{\c}{SplitTrain}{white}}
    \foreach \r/\c in {4/12,7/11} {\cella{\r}{\c}{SplitVal}{black}}
    \foreach \r/\c in {2/13,2/14,3/13,3/14,4/13,4/14,4/15,5/13,5/14,5/15,6/13,6/14,6/15,7/13}
      {\cella{\r}{\c}{SplitTest}{white}}
    \node[font=\small\bfseries, anchor=west] at (5.6,-0.3) {(a) standard};
  \end{scope}

  % ---------------- (b) tesi del 2022 ------------------------------------------
  \begin{scope}[xshift=76mm]
    \assi
    \foreach \r/\c in {2/10,2/11,2/12,2/13,2/14,3/10,6/7,6/13,7/13,7/9} {\esclusa{\r}{\c}}
    \foreach \r/\c in {3/11,3/12,3/13,4/10,4/12,4/13,4/14,5/10,5/11,5/12,5/14,5/15,
                       6/8,6/9,6/10,6/11,6/12,6/14,7/7,7/8,7/10,7/12}
      {\cella{\r}{\c}{SplitTrain}{white}}
    \foreach \r/\c in {3/14,4/15,5/13,6/15,7/11,4/11} {\cella{\r}{\c}{SplitVal}{black}}
    \node[font=\small\bfseries, anchor=west] at (5.6,-0.3) {(b) tesi del 2022};
  \end{scope}

  % ---------------- legenda --------------------------------------------------
  \begin{scope}[shift={(6.5,-8.4)}, font=\footnotesize]
    \draw[black!45, fill=SplitTrain] (0,-0.25) rectangle (0.5,0.25); \node[anchor=west] at (0.6,0) {training};
    \draw[black!45, fill=SplitVal] (3.2,-0.25) rectangle (3.7,0.25); \node[anchor=west] at (3.8,0) {validation};
    \draw[black!45, fill=SplitTest] (6.8,-0.25) rectangle (7.3,0.25); \node[anchor=west] at (7.4,0) {test};
    \draw[black!45, pattern=north east lines, pattern color=black!35] (9.5,-0.25) rectangle (10.0,0.25);
    \node[anchor=west] at (10.1,0) {esclusa};
  \end{scope}
\end{tikzpicture}
}
  \caption[Split delle tile di Potsdam]{Assegnazione delle 38 tile di Potsdam agli split.
    (a)~Protocollo standard della letteratura, usato per tutte le conclusioni. (b)~Protocollo
    della tesi del 2022: quattro delle aree escluse (\tile{2\_13}, \tile{2\_14}, \tile{6\_13},
    \tile{7\_13}) sono tile di test nel protocollo standard.}
  \label{fig:split}
\end{figure}

% ---------------------------------------------------------------------
\section{Preparazione dei dati}
\label{sec:preparazione}

Le maschere RGB vengono convertite in indici con una tabella di conversione vettorizzata sui
codici a 24 bit dei colori. I colori fuori palette diventano 255. Le tile di training vengono
tagliate in ritagli di $512\times512$ pixel con passo 512. L'ultimo ritaglio di ogni riga e
colonna è allineato al bordo, e i ritagli completamente non etichettati vengono scartati. Su
Potsdam si ottengono 144 ritagli per tile, quindi \num{3168} ritagli di training (22 tile) e 288 di
validation (2 tile). Su Vaihingen, con sole 14 aree di training, i ritagli si sovrappongono
(passo 256), ma solo all'interno del training. Le tile di validation e di test vengono registrate
anche intere, con la maschera piena e quella eroded, per la valutazione finale.

\paragraph{Regime a poche etichette.}
Per la \RQ{3} il sottoinsieme di ritagli etichettati viene campionato \emph{per tile} con seed
fisso. A parità di epoche, una frazione piccola di etichette produrrebbe pochissimi
aggiornamenti dei pesi (con l'1\% delle etichette, 2 batch per epoca) e il confronto tra
pre-addestramenti misurerebbe soprattutto il sotto-addestramento. I ritagli etichettati vengono
quindi ripetuti, con un nuovo rimescolamento a ogni passata, fino ad avere per epoca tanti batch
quanti con tutte le etichette. Ogni frazione riceve così lo stesso budget di iterazioni
(\cref{tab:frazioni}). Il pre-addestramento DINO usa comunque tutte le immagini di training,
senza etichette.

\begin{table}[t]
  \centering
  \small
  \caption[Regime a poche etichette su Potsdam]{Regime a poche etichette su Potsdam: ritagli
    etichettati e batch per epoca (batch effettivo 8, accumulo 4), con e senza ripetizione fino
    alla lunghezza di un'epoca completa.}
  \label{tab:frazioni}
  \begin{tabular}{|c|r|c|c|}
    \hline
    \textbf{Frazione} & \makecell{\textbf{Ritagli}\\\textbf{etichettati}} & \makecell{\textbf{Batch/epoca}\\\textbf{(con ripetizione)}} & \makecell{\textbf{Batch/epoca}\\\textbf{(senza)}} \\
    \hline
    100\% & \num{3168} & 396 & 396 \\
    25\%  & 792  & 396 & 99 \\
    10\%  & 308  & 396 & 38 \\
    1\%   & 22   & 396 & 2 \\
    \hline
  \end{tabular}
\end{table}

\paragraph{Pool non etichettati.}
Per verificare se più dati di dominio aiutino DINO, il pre-addestramento può usare anche i frame
grezzi di Toronto (RGB a 8 bit, \num{4433} ritagli) e di Vaihingen (IR-R-G a 16 bit, riportati a
8 bit con uno stretch ai percentili 0,5--99,5). I frame di Vaihingen coprono anche le aree di
test di Vaihingen, quindi un pool della stessa città del dataset di valutazione viene rifiutato
automaticamente, per evitare un'informazione trasduttiva sul test.

% ---------------------------------------------------------------------
\section{Configurazione di addestramento}
\label{sec:addestramento}

\subsection{Segmentazione}

La \cref{tab:iperparametri} riporta le configurazioni di addestramento. La configurazione di base
mantiene gli iperparametri del 2022 dove possibile (learning rate, StepLR, batch effettivo 32,
100 epoche) e corregge i difetti del protocollo. L'ottimizzatore è AdamW~\cite{loshchilov2019adamw}
(nel 2022 Adam~\cite{kingma2015adam}). Il batch di 8 ritagli con accumulo del gradiente su 4 passi
dà un batch effettivo di 32 nei limiti di memoria della GPU, con precisione mista
bfloat16~\cite{micikevicius2018mixed}. La normalizzazione con media e deviazione standard di
ImageNet è il primo strato del modello: training e inferenza non possono divergere su questo
punto. L'augmentation di training applica a immagine e maschera, in modo identico, riflessioni
orizzontali e verticali e rotazioni di \ang{90}: sono gli elementi di $D_4$, che non richiedono
interpolazione e lasciano le etichette esatte. Il miglior checkpoint viene scelto sulla mIoU di
validation, e l'addestramento si ferma dopo 20 epoche senza miglioramenti.

\begin{table}[t]
  \centering
  \small
  \caption[Configurazioni di addestramento della segmentazione]{Configurazioni di addestramento
    della segmentazione. La colonna \enquote{tesi 2022} riproduce il protocollo originale ed è
    usata solo per il confronto con la \cref{tab:risultati-2022}.}
  \label{tab:iperparametri}
  \begin{tabular}{|l|c|c|c|}
    \hline
    \textbf{Parametro} & \textbf{Base (ResNet)} & \textbf{Transformer} & \textbf{Tesi 2022} \\
    \hline
    Ritaglio / passo & 512 / 512 & 512 / 512 & 256 / 256 \\
    Augmentation & $D_4$ & $D_4$ & nessuna \\
    Ottimizzatore & AdamW & AdamW & Adam \\
    Learning rate & $10^{-3}$ & $10^{-4}$ & $10^{-3}$ \\
    Weight decay & $10^{-4}$ & $10^{-4}$ & -- \\
    Scheduler & StepLR (15; 0,4) & coseno & StepLR (15; 0,4) \\
    Batch $\times$ accumulo & $8\times4$ & $4\times8$ & $32\times1$ \\
    Epoche / early stopping & 100 / 20 & 100 / 20 & 100 / no \\
    Precisione & bf16 mista & bf16 mista & bf16 mista \\
    Dropout / drop path & 0,1 / -- & 0,1 / 0,2 & 0,1 / -- \\
    Valutazione: finestra / passo & 512 / 256 & 512 / 256 & 256 / 128 \\
    \hline
  \end{tabular}
\end{table}

\paragraph{Quantizzazione.}
La \maresunetvq{} usa un codebook di $K=512$ codici con $\beta=0{,}25$, aggiornamento EMA con
$\gamma=0{,}99$, inizializzazione dai dati e riavvio dei codici con conteggio EMA inferiore a
0,1. Durante l'addestramento si registrano perplexity e uso del codebook a ogni epoca.
L'ablation della \RQ{2} varia $K\in\{64,256,512,1024\}$, $\beta\in\{0{,}1;\,0{,}25;\,1{,}0\}$,
il peso $\lambda$, la quantizzazione delle connessioni di salto e l'aggiornamento per gradiente
($\gamma=0$) al posto dell'EMA.

\paragraph{Encoder transformer.}
Oltre a ResNet-18/34/50, l'encoder può essere uno Swin-T~\cite{liu2021swin} o un
ViT~\cite{dosovitskiy2021vit}. Il ViT ha una sola scala, e le quattro scale richieste dal decoder
vengono costruite con una \emph{simple feature pyramid}~\cite{li2022vitdet}. Tutti i pesi
ImageNet vengono portati da torchvision, con l'equivalenza numerica verificata da test.

\subsection{Pre-addestramento DINO}

La \cref{tab:iperparametri-dino} riporta la configurazione del pre-addestramento e la
\cref{tab:augmentation-dino} le due politiche di augmentation confrontate nella \RQ{4}.

\begin{table}[t]
  \centering
  \small
  \caption[Configurazione del pre-addestramento DINO]{Configurazione del pre-addestramento DINO,
    a confronto con quanto dichiarato e quanto implementato nel 2022.}
  \label{tab:iperparametri-dino}
  \begin{tabular}{|l|c|c|}
    \hline
    \textbf{Parametro} & \textbf{Questo lavoro} & \textbf{2022 (codice)} \\
    \hline
    Testa di proiezione & \makecell{MLP 2048-2048-256\\$K=\num{16384}$ prototipi} & \makecell{assente\\(uscita densa)} \\
    Normalizzazione della softmax & per campione & sull'intero batch \\
    Viste & $2\times224^2+6\times96^2$ & viste locali a $224^2$ \\
    Learning rate & \makecell{$5\cdot10^{-4}\cdot B/256$\\warmup 10 epoche, coseno} & fisso \\
    Weight decay & coseno $0{,}04\to0{,}4$ & assente \\
    $\tau_s$ / $\tau_t$ & \makecell{0,1 / $0{,}04\to0{,}07$\\in 30 epoche} & 0,1 / 0,04 fissa \\
    Momentum del teacher & coseno $0{,}996\to1$ & 0,996 fisso \\
    Teacher iniziale & copia dello student & rete indipendente \\
    Centering & $\lambda=0{,}9$, per campione & \makecell{$\lambda=0{,}9$, sul batch\\appiattito} \\
    Epoche / batch & 100 / 32 & 100 / -- \\
    Clipping del gradiente & norma 3,0 & -- \\
    \hline
  \end{tabular}
\end{table}

\begin{table}[t]
  \centering
  \small
  \caption[Politiche di augmentation per DINO]{Politiche di augmentation per DINO confrontate
    nella \RQ{4}. Le probabilità di blur sono 1,0 e 0,1 per le due viste globali e 0,5 per le
    locali; la solarizzazione si applica solo alla seconda vista globale.}
  \label{tab:augmentation-dino}
  \begin{tabular}{|l|c|c|}
    \hline
    \textbf{Trasformazione} & \textbf{Standard} & \textbf{Aerea} \\
    \hline
    Riflessione orizzontale & sì & sì \\
    Riflessione verticale, rotazioni \ang{90} & no & sì \\
    Jitter ($p=0{,}8$): luminosità, contrasto & 0,4; 0,4 & 0,2; 0,2 \\
    Jitter: saturazione, tinta & 0,2; 0,1 & 0,1; 0,02 \\
    Scala di grigi & $p=0{,}2$ & $p=0{,}1$ \\
    Solarizzazione (2\textsuperscript{a} vista globale) & $p=0{,}2$ & mai \\
    Blur gaussiano & sì & sì \\
    \hline
  \end{tabular}
\end{table}

% ---------------------------------------------------------------------
\section{Protocollo di valutazione}
\label{sec:valutazione}

La valutazione finale avviene sulle tile di test intere, con la finestra scorrevole della
\cref{sec:finestra} (finestra 512, passo 256) e i pesi del miglior checkpoint di validation. Ogni
tile viene predetta una volta sola e confrontata con due riferimenti, le etichette piene e quelle
eroded. Si accumulano così due matrici di confusione sull'intero split. Dalle matrici si
ricavano, per ciascun riferimento, IoU, F1, precision e recall per classe, OA, mF1 e mIoU, con
tutte le classi e senza clutter. Il formato riprende le tabelle 4.3--4.8 della tesi del 2022, per
facilitare il confronto. La metrica principale è la mIoU con tutte le classi su Potsdam e la mIoU
senza clutter su Vaihingen.

% ---------------------------------------------------------------------
\section{Implementazione}
\label{sec:implementazione}

Il codice del 2022 non è stato riusato. Il pacchetto \rsseg{} (\emph{Remote Sensing
SEGmentation}) è una reimplementazione da zero (\cref{fig:pipeline}): modelli e addestramento in
Keras~3~\cite{chollet2015keras} su backend PyTorch~\cite{paszke2019pytorch}, tracciamento degli
esperimenti con MLflow~\cite{zaharia2018mlflow}, configurazioni YAML componibili con validazione
delle chiavi. I parametri, la configurazione completa, il commit del codice, le metriche per
epoca, il picco di memoria della GPU, i checkpoint e i report di ogni esecuzione vengono
registrati su MLflow.

\begin{figure}[t]
  \centering
  \resizebox{\linewidth}{!}{% Pipeline sperimentale: dai dati ISPRS al report, con il ramo di pre-addestramento DINO
% e l'orchestrazione delle campagne. Coordinate esplicite su una griglia 5 x 2.
\begin{tikzpicture}[
    p/.style={attn, text width=2.4cm, minimum height=14mm},
    ds/.style={dato, text width=2.3cm, minimum height=12mm, font=\scriptsize},
  ]
  \def\dx{3.25}
  % riga superiore: segmentazione supervisionata
  \node[ds] (raw)   at (0,0)      {Dati ISPRS\\TOP, etichette, DSM/nDSM};
  \node[p]  (prep)  at (\dx,0)    {\texttt{prepare}\\{\scriptsize split per tile, LUT, ritagli, nDSM, statistiche}};
  \node[ds] (man)   at (2*\dx,0)  {ritagli e tile intere\\\texttt{crops.csv}, \texttt{split.json}};
  \node[p]  (train) at (3*\dx,0)  {\texttt{train}\\{\scriptsize Keras \texttt{fit}, CE + VQ, early stopping}};
  \node[p]  (eval)  at (4*\dx,0)  {valutazione\\{\scriptsize finestra scorrevole, etichette piene + eroded}};
  % riga inferiore: pre-addestramento
  \node[ds] (pool)  at (0,-2.6)     {frame non etichettati\\Toronto, Vaihingen};
  \node[p]  (unl)   at (\dx,-2.6)   {\texttt{prepare-}\\\texttt{unlabeled}\\{\scriptsize ritagli 512, 8 bit}};
  \node[p]  (dino)  at (2*\dx,-2.6) {\texttt{pretrain}\\{\scriptsize DINO, multi-crop}};
  \node[ds] (bb)    at (3*\dx,-2.6) {encoder dello student\\\texttt{.weights.h5}};
  \node[blocco, fill=SapRossoVelo, text width=2.4cm, minimum height=12mm] (mlf) at (4*\dx,-2.6)
       {MLflow\\{\scriptsize parametri, metriche, artefatti}};
  % orchestrazione
  \node[quant, text width=8.6cm] (camp) at (2.5*\dx,2.0)
       {\texttt{campaign}: passi in sottoprocessi, prerequisiti automatici, ripresa, report};

  \draw[flusso] (raw) -- (prep);
  \draw[flusso] (prep) -- (man);
  \draw[flusso] (man) -- (train);
  \draw[flusso] (train) -- (eval);
  \draw[flusso] (pool) -- (unl);
  \draw[flusso] (unl) -- (dino);
  \draw[flusso] (man) -- node[nota, right, pos=0.45] {immagini di training} (dino);
  \draw[flusso] (dino) -- (bb);
  \draw[flusso] (bb) -- node[nota, right, pos=0.45] {inizializzazione} (train);
  \draw[flusso] (eval) -- (mlf);
  \draw[flusso, densely dotted] (bb) -- (mlf);

  \draw[->, SapRossoScuro, densely dashed] (camp.west) -| (prep.north);
  \draw[->, SapRossoScuro, densely dashed] (camp.south -| train.north) -- (train.north);
  \draw[->, SapRossoScuro, densely dashed] (camp.east) -| (eval.north);
\end{tikzpicture}
}
  \caption[Pipeline sperimentale]{Pipeline sperimentale. La preparazione divide i dati per tile
    e produce ritagli e tile intere. Il ramo inferiore pre-addestra l'encoder con DINO, anche con
    pool non etichettati. Le campagne orchestrano i passi in sottoprocessi riprendibili, e ogni
    passo registra parametri, metriche e artefatti su MLflow.}
  \label{fig:pipeline}
\end{figure}

\paragraph{Verifica.}
109 test automatici, eseguiti in integrazione continua, verificano le proprietà da cui dipendono
le conclusioni (\cref{tab:verifiche}). Alcune sono equivalenze numeriche con riferimenti
indipendenti, altre sono invarianti dei dati, altre ancora test di regressione sui difetti del
2022.

\begin{table}[t]
  \centering
  \small
  \caption[Verifiche automatiche]{Principali verifiche automatiche dell'implementazione.}
  \label{tab:verifiche}
  \begin{tabular}{|l|p{9.6cm}|}
    \hline
    \textbf{Ambito} & \textbf{Proprietà verificata} \\
    \hline
    Encoder & ResNet-18/34 Keras = torchvision (pesi e statistiche BN casuali), Swin e ViT con padding e shift \\
    Attenzione & PAM lineare = trascrizione della formula originale = attenzione esplicita $1+\hatq\T\hatk$; CAM = riferimento NumPy \\
    VQ & gradiente straight-through = identità; perdita solo in training; l'EMA ritrova centri noti con perplexity alta \\
    DINO & perdita = media delle perdite per campione; teacher = copia iniziale, non addestrabile, aggiornato per EMA \\
    Dati & split disgiunti; conversione colori biunivoca; allineamento immagine/maschera dopo l'augmentation; frazioni di etichette deterministiche \\
    nDSM & encoder esteso con pesi nulli = encoder RGB (\cref{prop:inflazione}) \\
    Metriche & metriche accumulate su più batch = calcolo NumPy sulla stessa matrice di confusione \\
    Regressione & il clutter riceve gradiente (\cref{prop:classe-esclusa}) \\
    End-to-end & dati sintetici $\to$ preparazione $\to$ training $\to$ report $\to$ DINO $\to$ fine-tuning \\
    Ripresa & un training messo in pausa e ripreso coincide con uno ininterrotto \\
    \hline
  \end{tabular}
\end{table}

\paragraph{Campagne.}
Ogni domanda di ricerca è una \emph{campagna} YAML che elenca configurazioni, seed, sweep di
parametri, pre-addestramenti richiesti e benchmark. L'esecutore espande la campagna in passi con
un identificativo stabile e prepara i prerequisiti mancanti (dati, pool, pre-addestramenti).
Esegue poi ogni passo in un sottoprocesso, così che un errore o un esaurimento della memoria
fermi solo quel passo. I passi completati vengono riusati tra campagne diverse, e una campagna
interrotta riprende dal punto in cui si era fermata. Training e pre-addestramenti possono essere
messi in pausa e ripresi in modo esatto: l'ordine dei dati e le augmentation dipendono solo dal
seed e dalla posizione nel flusso.

% ---------------------------------------------------------------------
\section{Budget computazionale}
\label{sec:budget}

Tutti gli esperimenti sono dimensionati per una singola GPU NVIDIA RTX 4060 da \SI{8}{\giga\byte}.
La \cref{tab:vram} riporta il picco di memoria e il throughput misurati su passi di training
reali. ResNet-50 con batch 8 non entra in memoria e viene eseguito con batch 4 e accumulo 8. Il
tempo di un'epoca su Potsdam è di circa 2,8 minuti per \maresunet{} R18 e 3,2 minuti per R34. La
valutazione delle 14 tile di test richiede circa 7 minuti. L'insieme di tutte le campagne è
stimato in circa 16 giorni di calcolo, di cui circa 4 per il solo ViT-B/16.

\begin{table}[t]
  \centering
  \small
  \caption[Memoria e throughput misurati su RTX 4060]{Picco di memoria della GPU e throughput
    misurati su RTX 4060 (\SI{8}{\giga\byte}) con passi di training reali, ritagli $512^2$ e
    precisione bf16 mista.}
  \label{tab:vram}
  \begin{tabular}{|l|l|r|r|}
    \hline
    \textbf{Configurazione} & \textbf{Batch $\times$ accumulo} & \textbf{VRAM} & \textbf{Immagini/s} \\
    \hline
    \maresunet{} R18          & $8\times4$ & \SI{4,8}{\giga\byte} & 21 \\
    \maresunet{} R18 + nDSM   & $8\times4$ & \SI{4,8}{\giga\byte} & 23 \\
    \maresunetvq{} R18        & $8\times4$ & \SI{4,7}{\giga\byte} & 23 \\
    \maresunet{} R34          & $8\times4$ & \SI{5,5}{\giga\byte} & 19 \\
    \maresunet{} R50          & $8\times4$ & esaurita & -- \\
    \maresunet{} Swin-T       & $4\times8$ & \SI{4,2}{\giga\byte} & 13 \\
    \maresunet{} ViT-B/16     & $4\times8$ & \SI{7,1}{\giga\byte} & 2 \\
    DINO R18 ($2\times224^2+6\times96^2$) & $32$ & \SI{4,2}{\giga\byte} & 85 \\
    DINO R34                  & $32$ & \SI{6,2}{\giga\byte} & 40 \\
    \hline
  \end{tabular}
\end{table}

% ---------------------------------------------------------------------
\section{Analisi statistica e riproducibilità}
\label{sec:statistica}

Le differenze attese tra le varianti sono dell'ordine di uno o due punti di mIoU, comparabili con
la variabilità tra esecuzioni con seed diversi~\cite{bouthillier2021variance}. Per questo:

\begin{itemize}
  \item ogni configurazione viene addestrata con \textbf{tre seed} (0, 1, 2), e si riportano media
  e deviazione standard;
  \item i confronti sono \textbf{appaiati per seed}: il seed fissa l'inizializzazione delle parti
  non pre-addestrate, l'ordine dei dati e le augmentation, quindi la differenza
  $\Delta^{(s)}=\mathrm{mIoU}^{(s)}_{\mathrm{variante}}-\mathrm{mIoU}^{(s)}_{\mathrm{rif}}$
  elimina una parte della variabilità comune;
  \item la variabilità dovuta al campione di test (solo 14 tile su Potsdam) viene stimata con un
  \textbf{bootstrap sulle tile}~\cite{efron1993bootstrap}: si ricampionano con reinserimento le
  matrici di confusione delle singole tile e si ricalcola la mIoU, ottenendo un intervallo al
  95\%;
  \item un effetto viene considerato sostenuto dai dati solo se $\Delta^{(s)}$ ha lo stesso segno
  per tutti i seed \emph{e} l'intervallo bootstrap della differenza esclude lo zero; altrimenti
  il risultato è riportato come non conclusivo.
\end{itemize}

\Segnaposto{Implementare il bootstrap per tile: il report salva oggi la mIoU di ogni tile ma
non la sua matrice di confusione, che serve per ricalcolare la mIoU globale sui campioni
bootstrap. Fissare il numero di ricampionamenti.}

Ogni esecuzione salva la configurazione completa, il commit del codice e il manifest dei dati, e
gli split sono deterministici e verificati da test. I risultati di ogni tabella del
\cref{cap:esperimenti} si possono quindi rigenerare con un solo comando di campagna.
